\documentclass[12pt,a4paper]{article}
\usepackage[utf8]{inputenc}
\usepackage[spanish]{babel}
\usepackage{geometry}
\usepackage{amsmath}
\usepackage{amssymb}
\usepackage{booktabs}
\usepackage{longtable}
\usepackage{graphicx}
\usepackage{enumitem}
\usepackage{float}
\usepackage{xcolor}
\usepackage{fancyhdr}
\usepackage{caption}
\usepackage[numbers]{natbib}
\usepackage[colorlinks=true,linkcolor=black,citecolor=blue,urlcolor=blue]{hyperref}

\geometry{margin=2.5cm}

\pagestyle{fancy}
\fancyhf{}
\fancyhead[L]{Prueba Sistema Clasico}
\fancyhead[R]{\thepage}

\setcounter{tocdepth}{2}

\begin{document}

\begin{titlepage}
\centering
\vspace*{2cm}

{\Huge\bfseries Memoria de Cálculo\par}
\vspace{1.5cm}

{\Large\itshape Dimensionamiento de:\par}
\vspace{0.5cm}

{\LARGE\bfseries Prueba Sistema Clasico\par}
\vspace{2cm}

{\large Caudal de diseño: 17.0 L/s\par}
{\large Número de líneas: 3\par}
\vfill
{\large\today\par}
\end{titlepage}

\tableofcontents

\listoffigures

\listoftables

\newpage
\section{Resumen del Proyecto}

Este documento presenta el dimensionamiento detallado del Prueba Sistema Clasico, diseñado para tratar un caudal de 17.0 L/s distribuido en 3 líneas paralelas (5.67 L/s por línea).

\textbf{Características del afluente:}
\begin{itemize}
    \item DBO$_5$: 250.0 mg/L
    \item DQO: 500.0 mg/L
    \item SST: 220.0 mg/L
    \item Coliformes fecales: 1e+07 NMP/100mL
    \item Temperatura: 18.0°C
\end{itemize}

\textbf{Secuencia de unidades del tren:} 
Humedal Artificial de Flujo Vertical (HAFV)

\section{Dimensionamiento de Unidades}

\newpage

\subsection{Humedal Artificial de Flujo Vertical (HAFV)}

\subsubsection{Dimensionamiento}

\subsubsection*{Principios de Funcionamiento del Humedal Artificial de Flujo Vertical}

El Humedal Artificial de Flujo Vertical (HAFV) opera mediante percolación gravitacional del agua residual a través de un medio filtrante granular estratificado, donde se desarrolla biopelícula aerobia responsable de la degradación de materia orgánica y nitrificación. La alimentación intermitente genera ciclos de saturación y drenaje que re-airean pasivamente el lecho, favoreciendo la estabilización de sólidos retenidos.

Durante la fase de alimentación, el agua residual se distribuye sobre la superficie del filtro mediante un sistema de tuberías perforadas y percola hacia las capas inferiores, donde es recolectada por las tuberías de drenaje. Durante el período de reposo subsiguiente, el medio drena libremente y el aire penetra por convección natural desde los tubos de ventilación del fondo y desde la superficie expuesta, restaurando las condiciones aerobias en el interior del lecho. Este ciclo de saturación-drenaje-re-aireación, además de proveer el oxígeno necesario para las reacciones biológicas, permite la descomposición aerobia progresiva de los sólidos acumulados en la capa superficial del filtro, retardando la colmatación y extendiendo la vida útil del sistema sin necesidad de intervenciones frecuentes.

Los mecanismos de remoción y las eficiencias típicas del HAFV se presentan en la Tabla~\ref{tab:mecanismos_hafv_a}.

\begin{table}[H]
\centering
\caption{Mecanismos de remoción y eficiencias típicas en el HAFV -- Sistema Clásico (Kadlec \& Wallace, 2009; Cooper et al., 1996)}
\label{tab:mecanismos_hafv_a}
\begin{tabular}{p{2.2cm}p{6.5cm}cc}
\toprule
\textbf{Parámetro} & \textbf{Mecanismo principal} & \textbf{Eficiencia típica} & \textbf{Unidad} \\
\midrule
DBO$_5$ / DQO & Degradación aerobia y anóxica en biopelícula & 70 -- 85 & \% \\
SST & Filtración física en el medio granular & 80 -- 90 & \% \\
NH$_4$-N & Nitrificación aerobia (\textit{Nitrosomonas/Nitrobacter}) & 40 -- 60 & \% \\
Colif. fecales & Filtración, adsorción, oxidación biológica & 1 -- 2 & log \\
Fósforo total & Adsorción en el medio (limitada) & 20 -- 40 & \% \\
\bottomrule
\end{tabular}
\end{table}

\subsubsection*{Criterio de Selección del Sistema: Sistema Clásico en Clima Templado-Frío}

La selección entre las dos metodologías de diseño disponibles para el HAFV se fundamenta en la temperatura del agua residual como criterio primario. A temperaturas inferiores a 15\,°C, la cinética de mineralización aerobia del lodo superficial es insuficiente para sostener ciclos de reposo cortos sin riesgo de colmatación prematura. En consecuencia, el diseño debe ser conservador en cargas y generoso en períodos de reposo.

Para la temperatura de \textbf{18.0\,°C} del presente proyecto, inferior al umbral de 20\,°C, corresponde aplicar el \textbf{Sistema Clásico (Ruta A)}, conforme a la metodología de Cooper et al. (1996) y la norma austriaca ÖNORM B 2505 (2009). A esta temperatura, la actividad microbiana es más lenta y la mineralización del lodo superficial requiere períodos de reposo más prolongados (7 días o más por filtro). La configuración de dos etapas en serie distribuye la carga de forma decreciente: la primera etapa, más cargada, remueve la fracción principal de la DBO$_5$ soluble y suspendida; la segunda etapa actúa como unidad de pulido, tratando el efluente de la primera con cargas menores y produciendo una calidad superior a la que podría alcanzar una etapa única.

\subsubsection*{Configuración en Dos Etapas y Ciclo de Operación}

La disposición en dos etapas en serie es característica del Sistema Clásico. La primera etapa opera con 3 filtros en paralelo bajo un ciclo de 3.5 días de alimentación por 7.0 días de reposo, con lo cual en todo momento hay un solo filtro recibiendo caudal mientras los otros dos se recuperan. La segunda etapa opera con 2 filtros en paralelo bajo ciclos de 7 días de alimentación por 7 días de reposo. Esta diferencia en los ciclos refleja la mayor carga de la primera etapa y la necesidad de períodos de recuperación más extensos a menor temperatura.

El medio filtrante del Sistema Clásico difiere del sistema francés en su capa superficial: mientras el sistema francés emplea grava fina (2 -- 6\,mm), el clásico utiliza arena gruesa (0 -- 4\,mm) en la capa superior, seguida de grava media (4 -- 8\,mm) en la capa principal y grava gruesa (16 -- 32\,mm) en la capa de drenaje. La arena superficial proporciona mayor área específica de contacto con el agua y favorece la retención de sólidos, pero requiere ciclos de reposo largos para su mineralización a temperaturas bajas.

\subsubsection*{Dimensionamiento por Carga Orgánica Superficial (COS)}

La Carga Orgánica Superficial (COS) expresa la masa de DBO$_5$ aplicada por unidad de área del filtro activo por día. Es el criterio de dimensionamiento principal del Sistema Clásico, dado que a bajas temperaturas el control biológico del proceso es el factor limitante. El valor adoptado de 40\,g\,DBO$_5$/m²·d corresponde al rango recomendado para la primera etapa en clima templado (Cooper et al., 1996), y garantiza que la biopelícula dispone de suficiente tiempo de contacto para degradar la carga recibida incluso bajo condiciones de menor actividad metabólica. El área de la primera etapa por este criterio es:

\begin{equation}
A_{1a} = \frac{Q \times S_{0,\text{DBO5}} \times 10^{-3}}{\text{COS}_1} = \frac{489.6 \times 75 \times 10^{-3}}{40 \times 10^{-3}} = 1224.0\ \text{m}^2
\end{equation}
\captionequation{Área de la primera etapa por criterio de carga orgánica superficial (COS) -- HAFV Sistema Clásico}

La segunda etapa recibe el efluente de la primera con carga ya parcialmente removida. Su área de diseño se establece en el 50\,\% del área de la primera etapa, lo que refleja la menor carga orgánica recibida y los menores requerimientos de tiempo de contacto para el pulido final:

\begin{equation}
A_{2a} = 0.50 \times A_{1a} = 0.50 \times 1224.0 = 612.0\ \text{m}^2
\end{equation}
\captionequation{Área de la segunda etapa (50\,\% de la primera etapa) -- HAFV Sistema Clásico}

El área total del sistema es la suma de ambas etapas:

\begin{equation}
A_\text{total} = A_{1a} + A_{2a} = 1224.0 + 612.0 = 1836.0\ \text{m}^2
\end{equation}
\captionequation{Área total del sistema HAFV Sistema Clásico (suma de etapas)}

\subsubsection*{Verificación por Habitante Equivalente (PE)}

El criterio del habitante equivalente constituye una verificación empírica independiente del dimensionamiento. La unidad PE se define como 60\,g\,DBO$_5$/hab·d (estándar europeo), y permite relacionar la carga real del proyecto con referencias de diseño derivadas de la experiencia en instalaciones similares. En sistemas donde el HAFV recibe efluente de un proceso de tratamiento primario, el PE debe calcularse a partir de la carga real de DBO$_5$ al HAFV y no del caudal per cápita, para no sobredimensionar el sistema:

\begin{equation}
\text{PE}_\text{carga} = \frac{Q \times S_{0,\text{DBO5}}}{60 \times 10^{-3}} = \frac{489.6 \times 75 \times 10^{-3}}{60 \times 10^{-3}} = 612\ \text{hab. eq.}
\end{equation}
\captionequation{Habitante equivalente calculado a partir de la carga real de DBO$_5$ en el afluente al HAFV}

Con un área de referencia de 2.0\,m²/PE para la primera etapa en clima templado (Cooper et al., 1996), el área total del sistema por este criterio resulta:

\begin{equation}
A_\text{PE} = 1.50 \times a_1 \times \text{PE}_\text{carga} = 1.50 \times 2.0 \times 612 = 1836.0\ \text{m}^2
\end{equation}
\captionequation{Área de referencia por criterio de habitante equivalente -- Sistema Clásico}

no se aplicó el criterio de PE como controlante El área definitiva de diseño adoptada es $A_\text{total} = \mathbf{1836.0\ \text{m}^2}$.

\subsubsection*{Geometría de los Filtros}

Las unidades de la primera etapa presentan dimensiones de \textbf{28.6\,m $\times$ 14.3\,m} por filtro (relación largo:ancho de 2.0:1), con una profundidad de medio filtrante de 0.60 a 0.80\,m conforme a lo recomendado por Cooper et al. (1996) para la primera etapa del sistema clásico. La altura libre sobre el medio (0.50\,m) permite acomodar la acumulación de lodo superficial durante la vida útil estimada de 10 a 15 años.


\subsubsection{Verificación}

\subsubsection*{Fundamento del Modelo Cinético k-C*}

El dimensionamiento por carga orgánica e hidráulica superficial determina el área del sistema a partir de criterios empíricos derivados de la experiencia en instalaciones reales. Sin embargo, este enfoque no garantiza por sí solo que el sistema producirá la calidad de efluente requerida, ya que no incorpora explícitamente la cinética del proceso biológico ni las condiciones específicas de temperatura del proyecto. Por esta razón, toda propuesta de HAFV debe complementarse con una verificación de eficiencia mediante el modelo cinético de primer orden con concentración de fondo, conocido como modelo k-C* (Kadlec \& Wallace, 2009).

El modelo k-C* es el estándar de referencia internacional para la verificación de eficiencia en humedales construidos, adoptado por la guía EPA para humedales de tratamiento (US~EPA, 2000) y ampliamente validado en instalaciones de clima templado y tropical. Su formulación matematiza dos observaciones fundamentales de la práctica en humedales: en primer lugar, que la remoción de contaminantes sigue una cinética de primer orden en cuanto a la concentración activa del contaminante (fracción por encima de la concentración de fondo); y en segundo lugar, que existe una concentración mínima irreducible, denominada concentración de fondo o \textit{background concentration} ($C^*$), por debajo de la cual el sistema no puede reducir la concentración del efluente independientemente del área disponible. Esta concentración de fondo refleja la DBO$_5$ soluble de naturaleza recalcitrante, la demanda endógena generada por la propia biomasa del sistema y la resuspensión de material orgánico particulado.

La ecuación del modelo, en su forma para un reactor de flujo pistón (aproximación válida para HAFV de geometría rectangular), es:

\begin{equation}
C_e = C^* + (C_i - C^*) \cdot \exp\!\left(\frac{-k_{A,T}}{q}\right)
\label{eq:kC_star}
\end{equation}
\captionequation{Modelo cinético k-C* para humedales de flujo vertical -- Kadlec \& Wallace (2009)}

donde los términos de la ecuación tienen el siguiente significado físico:

\begin{itemize}[noitemsep,leftmargin=2em]
    \item[$C_e$] = Concentración de DBO$_5$ en el efluente calculada por el modelo [mg/L]
    \item[$C_i$] = Concentración de DBO$_5$ en el afluente al HAFV (75.0\,mg/L en este proyecto)
    \item[$C^*$] = Concentración de fondo (\textit{background concentration}): mínimo inalcanzable del sistema, fijado en 3.5\,mg\,DBO$_5$/L según Kadlec \& Wallace (2009) para HAFV con afluente municipal
    \item[$k_{A,T}$] = Constante cinética de remoción de primer orden a la temperatura $T$ del proyecto [m/d], que cuantifica la capacidad de remoción de la biopelícula por unidad de área y de carga hidráulica
    \item[$q$] = Carga hidráulica superficial del filtro activo [m/d], igual a $Q / A_\text{operando}$
\end{itemize}

El cociente $k_{A,T}/q$ es el argumento del exponencial y representa el número de unidades cinéticas de remoción disponibles: a mayor área (menor $q$) o mayor temperatura (mayor $k_{A,T}$), la remoción es más intensa. Cuando $k_{A,T}/q \gg 1$, la concentración de salida tiende asintóticamente hacia $C^*$; cuando $k_{A,T}/q \ll 1$, la remoción es mínima.

\subsubsection*{Corrección de la Constante Cinética por Temperatura}

La constante cinética $k_A$ depende fuertemente de la temperatura del agua residual, dado que la velocidad de las reacciones enzimáticas en la biopelícula es una función exponencial de la temperatura. Kadlec \& Wallace (2009) proponen la corrección mediante el modelo de Arrhenius modificado:

\begin{equation}
k_{A,T} = k_{A,20} \cdot \theta^{(T-20)}
\label{eq:arrhenius}
\end{equation}
\captionequation{Corrección de la constante cinética por temperatura -- Ecuación de Arrhenius modificada (Kadlec \& Wallace, 2009)}

En esta expresión, $k_{A,20}$ es la constante cinética de referencia a 20\,°C, determinada estadísticamente a partir de datos de seguimiento de numerosas instalaciones reales. Para la remoción de DBO$_5$ en humedales de flujo vertical tratando aguas residuales municipales, Kadlec \& Wallace (2009) establecen $k_{A,20} = 0.093$\,m/d. El coeficiente de temperatura $\theta = 1.06$ cuantifica la sensibilidad de la reacción a los cambios térmicos: por cada grado centígrado de aumento en la temperatura del agua, la constante cinética aumenta aproximadamente un 6\,\%.

La Tabla~\ref{tab:kAT_temperatura} ilustra la variación de la constante cinética con la temperatura para el parámetro DBO$_5$, permitiendo apreciar el efecto cuantitativo de la temperatura en el rendimiento del sistema.

\begin{table}[H]
\centering
\caption{Valores de $k_{A,T}$ para DBO$_5$ en función de la temperatura ($k_{A,20} = 0.093$\,m/d; $\theta = 1.06$)}
\label{tab:kAT_temperatura}
\begin{tabular}{ccc}
\toprule
\textbf{Temperatura (°C)} & \textbf{Factor $\theta^{(T-20)}$} & \textbf{$k_{A,T}$ (m/d)} \\
\midrule
10 & 0.558 & 0.052 \\
15 & 0.747 & 0.069 \\
18 & 0.890 & 0.083 \\
20 & 1.000 & 0.093 \\
22 & 1.124 & 0.105 \\
25 & 1.338 & 0.124 \\
18.0 & 1.06$^{(18.0-20)}$ & \textbf{0.0828} \\
28 & 1.594 & 0.148 \\
30 & 1.791 & 0.167 \\
\bottomrule
\end{tabular}
\end{table}

Para la temperatura de diseño del proyecto ($T = 18.0$\,°C), la constante cinética corregida resulta:

\begin{equation}
k_{A,T} = 0.093 \times 1.06^{(18.0 - 20)} = 0.093 \times 1.06^{18.0 - 20:.1f} = 0.0828\ \text{m/d}
\end{equation}
\captionequation{Constante cinética de remoción de DBO$_5$ corregida a $T = 18.0$\,°C}

\subsubsection*{Aplicación del Modelo y Estimación del Efluente}

La carga hidráulica superficial sobre el filtro activo, considerando únicamente el área en operación simultánea ($A_\text{operando}$), es:

\begin{equation}
q = \frac{Q}{A_\text{operando}} = \frac{489.6}{1224.0} = 0.4000\ \text{m/d}
\end{equation}
\captionequation{Carga hidráulica superficial del filtro activo ($q$) -- parámetro de entrada del modelo k-C*}

Sustituyendo en la ecuación del modelo~\eqref{eq:kC_star} con los valores del proyecto:

\begin{equation}
C_e = 3.5 + (75.0 - 3.5) \cdot \exp\!\left(\frac{-0.0828}{0.4000}\right) = 3.5 + 71.5 \cdot \exp(-0.2070) = 61.6\ \text{mg/L}
\end{equation}
\captionequation{Concentración estimada de DBO$_5$ en el efluente del HAFV -- Modelo k-C*}

La eficiencia de remoción de DBO$_5$ en el HAFV resulta:

\begin{equation}
\eta = \frac{C_i - C_e}{C_i} \times 100 = \frac{75.0 - 61.6}{75.0} \times 100 = 17.8\ \%
\end{equation}
\captionequation{Eficiencia de remoción de DBO$_5$ en el HAFV}

La Tabla~\ref{tab:verificacion_kC} consolida los parámetros y resultados de la verificación cinética.

\begin{table}[H]
\centering
\caption{Resultados de la verificación cinética k-C* -- Humedal Artificial de Flujo Vertical}
\label{tab:verificacion_kC}
\begin{tabular}{lccc}
\toprule
\textbf{Parámetro} & \textbf{Símbolo} & \textbf{Valor} & \textbf{Unidad} \\
\midrule
Constante cinética de referencia a 20\,°C & $k_{A,20}$ & 0.093 & m/d \\
Coeficiente de temperatura & $\theta$ & 1.06 & -- \\
Temperatura del agua residual & $T$ & 18.0 & °C \\
Constante cinética corregida por temperatura & $k_{A,T}$ & 0.0828 & m/d \\
Concentración de fondo & $C^*$ & 3.5 & mg\,DBO$_5$/L \\
Carga hidráulica del filtro activo & $q$ & 0.4000 & m/d \\
Relación $k_{A,T} / q$ & -- & 0.2070 & -- \\
DBO$_5$ del afluente & $C_i$ & 75.0 & mg/L \\
DBO$_5$ del efluente estimada (modelo k-C*) & $C_e$ & 61.6 & mg/L \\
DBO$_5$ objetivo de diseño & -- & 20.0 & mg/L \\
Eficiencia de remoción de DBO$_5$ en HAFV & $\eta$ & 17.8 & \% \\
\bottomrule
\end{tabular}
\end{table}

\textbf{Resultado de la verificación:} La DBO5 calculada de 61.6 mg/L no cumple con el objetivo de diseño de 20.0 mg/L.

\subsubsection*{Retención de Sólidos, Producción de Lodo y Mantenimiento}

El HAFV retiene los sólidos suspendidos del afluente mediante filtración física en la capa superficial del medio granular. Este material, denominado lodo superficial, se acumula progresivamente sobre la grava de la capa superior durante los ciclos de alimentación. La clave del sistema para evitar la colmatación radica en la mineralización aerobia de este lodo durante los períodos de reposo: a temperatura del agua residual de 18.0\,°C, la actividad microbiana degrada la fracción orgánica del lodo acumulado, reduciendo su volumen y manteniendo la permeabilidad del medio.

La producción neta de sólidos retenidos puede estimarse como:

\begin{equation}
P_\text{sólidos} = Q \times (S_{0,\text{SST}} - C_{e,\text{SST}}) \times 10^{-3}\ \text{[kg\,SST/d]}
\end{equation}
\captionequation{Producción de sólidos retenidos en el medio filtrante del HAFV}

En condiciones tropicales ($T > 20$\,°C), la acumulación neta de lodo superficial, una vez descontada la mineralización durante el reposo, es del orden de 0.5 a 1.0\,cm/año (Molle et al., 2015). Con una altura libre sobre el medio de 0.55\,m, el sistema puede operar sin retiro de lodo por períodos superiores a 10 años antes de requerir una evaluación de colmatación. El lodo eventualmente extraído durante el mantenimiento mayor tiene carácter de biosólido estabilizado aerobiamente y puede recibir disposición controlada sin requerir tratamiento adicional.

El programa de mantenimiento preventivo incluye inspección diaria de los sistemas de distribución y control de la alternancia de filtros, revisión semanal del drenaje completo durante el período de reposo, muestreo mensual del efluente para DBO$_5$, SST y coliformes fecales, cosecha semestral parcial de la vegetación macrófita (no más del 50\,\% de la biomasa aérea para no estresar las plantas), e inspección anual del espesor de lodo acumulado sobre el medio filtrante.


\subsubsection{Resultados}

La Tabla~\ref{tab:resumen_hafv} concentra los parámetros geométricos, hidráulicos, de carga y de calidad del efluente del sistema dimensionado, permitiendo una lectura integrada de los resultados del diseño y facilitando su verificación cruzada con los criterios de diseño aplicados.

\begingroup
\small
\setlength{\tabcolsep}{4pt}
\renewcommand{\arraystretch}{1.12}
\begin{longtable}{p{4.1cm}p{4.8cm}p{4.2cm}}
\caption{Resumen de resultados del dimensionamiento -- Humedal Artificial de Flujo Vertical (HAFV)}\label{tab:resumen_hafv}\\
\toprule
\textbf{Parámetro} & \textbf{Descripción} & \textbf{Valor} \\
\midrule
\endfirsthead
\caption[]{Resumen de resultados del dimensionamiento -- Humedal Artificial de Flujo Vertical (HAFV) (continuación)}\\
\toprule
\textbf{Parámetro} & \textbf{Descripción} & \textbf{Valor} \\
\midrule
\endhead
\midrule
\multicolumn{3}{r}{\textit{Continúa en la siguiente página}}\\
\endfoot
\bottomrule
\endlastfoot
\multicolumn{3}{l}{\textit{Metodología y configuración}} \\
\quad Sistema adoptado & Metodología & Clásico (Ruta A) \\
\quad Criterio controlante & Factor limitante & COS/PE \\
\quad N° celdas (alternancia) & Unidades en operación alterna & 3 \\
\quad Ciclo alimentación / reposo & Rotación por unidad & 3.5\,/\,7.0\,días \\
\midrule
\multicolumn{3}{l}{\textit{Geometría del sistema}} \\
\quad Área total por línea & Superficie celdas por línea & 1836.0\,m² \\
\quad Área total sistema & Superficie celdas (3 líneas) & 5508.0\,m² \\
\quad Área celda activa & Superficie celda en operación & 1224.0\,m² \\
\quad Dimensiones por celda & Largo $\times$ Ancho & 28.6$\times$14.3\,m \\
\quad Profundidad del medio & Espesor lecho filtrante & 0.70\,m \\
\quad Altura total & Incluye borde libre & 1.20\,m \\
\midrule
\multicolumn{3}{l}{\textit{Parámetros de carga}} \\
\quad OLR real (celda activa) & Carga orgánica superficial & 30.0\,g DQO/m²·d \\
\quad HLR real (celda activa) & Carga hidráulica superficial & 0.400\,m/d \\
\quad PE equivalente & Habitantes equivalentes & 612\,PE \\
\midrule
\multicolumn{3}{l}{\textit{Calidad del efluente estimada}} \\
\quad DBO$_5$ calculada (k-C*) & Concentración estimada & 61.6\,mg/L \\
\end{longtable}
\endgroup

El área total por línea de 1836.0\,m² resulta del criterio de diseño dominante (\textbf{COS/PE}), lo que implica que el factor limitante del sistema es de naturaleza orgánica: la capacidad de degradación de la biopelícula es el factor que determina el área, y la carga hidráulica resultante se sitúa dentro del rango admisible para el sistema seleccionado. En conjunto, los parámetros de carga reales verifican el cumplimiento de los rangos validados para el sistema Clásico (Molle et al., 2015).

La concentración de DBO$_5$ estimada en el efluente mediante el modelo k-C* es 61.6\,mg/L. Es importante señalar que la eficiencia de remoción de DBO$_5$ en el propio HAFV puede parecer moderada en sistemas con tratamiento primario previo, dado que el reactor anaerobio ya ha removido la fracción principal de la carga orgánica; la eficiencia total del tren de tratamiento es sustancialmente mayor.

Con respecto a la remoción de sólidos suspendidos, el HAFV presenta eficiencias típicas del 80 al 95\,\% por filtración física en el medio granular. Los coliformes fecales se reducen entre 1 y 3 unidades logarítmicas; si es necesario reducir más la concentración de coliformes fecales, se deberá incorporar una etapa de desinfección posterior (radiación ultravioleta o hipoclorito de sodio).

La Figura~\ref{fig:humedal_vertical} presenta el esquema del sistema con la vista en planta del arreglo de filtros en operación alternada y la sección transversal simplificada que ilustra la composición estratificada del medio filtrante y la ubicación del sistema de drenaje inferior.

\begin{figure}[H]
\centering
\includegraphics[width=0.85\textwidth]{C:/Users/Alienware/OneDrive/00_OFC/12_documentos_ofc/01_ToR/01_san_cristobal/00_ptar_cristobal/00_codigos/resultados/figuras/Esquema_Linea_Humedal_5Ls_2celdas.png}
\caption{Esquema del Humedal Artificial de Flujo Vertical -- Clásico. \textit{Arriba:} vista en planta con los 3 filtros en operación alternada (alimentando / reposo). \textit{Abajo:} sección transversal simplificada con las capas del medio filtrante (grava fina superior, grava media intermedia y grava gruesa de drenaje) y el sistema de tuberías de drenaje inferior con tubos de ventilación.}
\label{fig:humedal_vertical}
\end{figure}

\subsubsection*{Ciclo de Operación y Esquema de Rotación}

El esquema de operación alternada asegura que, mientras 2 unidad(es) permanece(n) en período de reposo durante 7.0 días, el filtro activo recibe la totalidad del caudal de diseño durante 3.5 días. Al término del ciclo de alimentación, los roles se invierten automáticamente mediante temporizador o sistema básico de control. Esta rotación distribuye el desgaste biológico entre las unidades, permite la recuperación aerobia completa del medio filtrante en cada ciclo y minimiza el riesgo de colmatación progresiva.

Durante el período de reposo, el drenaje completo del lecho es condición necesaria e indispensable: solo con el medio drenado pueden penetrar el oxígeno y el aire desde la superficie libre y los tubos de ventilación del fondo, habilitando la mineralización aerobia del lodo superficial. Cualquier obstrucción de los tubos de drenaje que impida el vaciado completo del sistema compromete directamente la eficiencia del proceso de regeneración y, a mediano plazo, conduce a la colmatación del filtro.

\subsubsection*{Vegetación Macrófita: Establecimiento y Manejo}

La vegetación macrófita (\textit{Phragmites australis} o especies tropicales equivalentes nativas) se establece sobre el lecho filtrante antes de la puesta en marcha del sistema, a una densidad de 4 a 6 plantas por metro cuadrado. Su sistema radicular, a medida que se desarrolla durante los primeros meses de operación, cumple funciones complementarias que contribuyen a la estabilidad del sistema:

\begin{itemize}[noitemsep,leftmargin=2em]
    \item Mantiene la permeabilidad superficial del medio filtrante al impedir la compactación de las capas superiores
    \item Crea micrositios aerobios localizados en la rizosfera, ampliando la diversidad de condiciones redox disponibles para las comunidades microbianas
    \item Transfiere pequeñas cantidades de oxígeno hacia las zonas profundas del lecho a través de los aerénquimas radiculares
    \item Contribuye a la integración paisajística del sistema en el entorno natural del proyecto
\end{itemize}

Durante los primeros 6 a 12 meses de operación (período de establecimiento de la biopelícula y de la vegetación), se recomienda operar el sistema con una carga reducida al 30 a 50\,\% de la carga de diseño, incrementándola gradualmente hasta alcanzar las condiciones nominales. Este período de arranque permite el desarrollo de la biopelícula, el acondicionamiento del medio filtrante y la estabilización de la vegetación antes de enfrentar la carga plena.



\newpage
\section{Referencias Bibliograficas}
\bibliographystyle{plain}
\bibliography{referencias}

\end{document}